\section*{Chapter \ref{ch:b3:heterocycles} --- Heterocyclic Chemistry}

\begin{solution}{exo:b3:heterocycles:1}
Oxazole: O gives two (one lone pair in the $\pi$ system, one in the plane), N one, the
three carbons three: 6. Thiazole: the same with S: 6. Pyrimidine: one from each
of the six atoms, both nitrogen lone pairs in the plane: 6. Purine: nine atoms,
the N--H nitrogen giving two and the others one each: 10, a Hückel number
($4n + 2$, $n = 2$).
\end{solution}

\begin{solution}{exo:b3:heterocycles:2}
Piperidine (11.1) $>$ DMAP (9.6) $>$ imidazole (7.0) $>$ pyridine (5.2) $\gg$ pyrrole (about $-4$,
protonated on carbon).
\end{solution}

\begin{solution}{exo:b3:heterocycles:3}
2-Bromopyrrole (pyrrole is so reactive that it is easily polybrominated);
2-bromofuran; 3-bromoindole.
\end{solution}

\begin{solution}{exo:b3:heterocycles:4}
2,5-Dimethylfuran; 1,2,5-trimethylpyrrole; 2,5-dimethylthiophene.
\end{solution}

\begin{solution}{exo:b3:heterocycles:5}
2-Methoxypyridine. Addition of methoxide at C2 puts the negative charge of the
\omterm{def:b3:heterocycles:snar}{Meisenheimer complex} on nitrogen; at C3 it can only be spread over carbons, so the
intermediate, and the transition state before it, are much higher in energy.
\end{solution}

\begin{solution}{exo:b3:heterocycles:6}
2-Aminopyridine (as its sodium salt, then protonated on work-up). The amide ion
adds at C2, giving an anionic $\sigma$-adduct with H and $\ce{NH2}$ on C2 and the negative
charge on the ring nitrogen; it loses hydride, which takes a proton from the
amino group and leaves as $\ce{H2}$.
\end{solution}

\begin{solution}{exo:b3:heterocycles:7}
The oxygen of the N-oxide gives a lone pair back into the ring: resonance
structures carry negative charge at C2, C4 and C6. Electrophiles attack C4 (C2
is close to the positively charged nitrogen). The oxygen is then removed with
phosphorus trichloride, giving 4-nitropyridine.
\end{solution}

\begin{solution}{exo:b3:heterocycles:8}
2-Hydroxypyridine (an OH on C2) and 2-pyridone (N--H and C=O). The pyridone keeps six
$\pi$ electrons in the ring: its zwitterionic resonance structure, with $\mathrm{O^-}$ and an
$\mathrm{N^+}$ that shares its lone pair, is a pyridinium-olate; it is an amide whose
nitrogen lone pair completes the sextet, and the strong C=O and N--H
hydrogen bonds favour it in water.
\end{solution}

\begin{solution}{exo:b3:heterocycles:9}
2-Nitrobenzaldehyde, two equivalents of methyl acetoacetate, and ammonia: the
calcium-channel blocker nifedipine.
\end{solution}

\begin{solution}{exo:b3:heterocycles:10}
The ene-hydrazine towards the $\ce{CH2}$ group (internal, more substituted) gives
2,3-dimethylindole; that towards the methyl group (terminal) gives 2-ethylindole.
Weak acids give mainly the first, through the more stable ene-hydrazine;
stronger acids raise the proportion of 2-ethylindole.
\end{solution}

\begin{solution}{exo:b3:heterocycles:11}
In $\mathrm S_\mathrm NAr$ the rate-determining step is the addition, which does not break the C--X
bond; the more electronegative fluorine stabilises the anionic transition state
and the \omterm{def:b3:heterocycles:snar}{Meisenheimer complex} more. In $\mathrm S_\mathrm N2$ the C--X bond breaks in the
transition state, and the weak C--I bond leaves fastest.
\end{solution}

\begin{solution}{exo:b3:heterocycles:12}
Pyridine: three signals in the ratio 2\,:\,1\,:\,2, the first near \qty{8.6}{ppm}. Pyrrole: two equal
signals at 6.7 and \qty{6.2}{ppm} and a broad N--H near \qty{8.1}{ppm}. Furan: two equal signals at
7.4 and \qty{6.4}{ppm}, with no N--H. Thiophene: two equal signals at 7.33 and \qty{7.12}{ppm}, close
together, between those of furan.
\end{solution}

\begin{solution}{pb:b3:heterocycles:1}
\textbf{1.} $\ce{C6H5NHNH2 + (CH3)2CO -> C6H5NHN=C(CH3)2 + H2O}$.

\textbf{2.} The amine nitrogen adds to the carbonyl carbon, then water is lost, as in
imine formation; acid activates the carbonyl group and turns the OH of the adduct
into a leaving group.

\textbf{3.} The terminal $\ce{NH2}$: the other nitrogen's lone pair is conjugated with the
phenyl ring and more hindered.

\textbf{4.} $\qty{108.0}{g/mol}$; $\qty{10.0}{g}/\qty{108.0}{g/mol} = \qty{0.0926}{mol}$.

\textbf{5.} The hydrazone $\ce{C9H12N2}$, \qty{148.0}{g/mol}: $\qty{0.0926}{mol} \times \qty{148.0}{g/mol} = \qty{13.7}{g}$.

\textbf{6.} $\ce{C6H5-NH-NH-C(CH3)=CH2}$.

\textbf{7.} The \textit{ortho} ring carbon, the \textit{ipso} carbon, the two nitrogens, the hydrazone carbon
and the terminal $\ce{CH2}$.

\textbf{8.} The N--N bond breaks; a C--C bond forms between the \textit{ortho} carbon and the $\ce{CH2}$.

\textbf{9.} Six electrons, all components \omterm{def:b3:pericyclic:faciality}{suprafacial}: thermally allowed.

\textbf{10.} As a Lewis acid it catalyses the tautomerisation, binds a nitrogen and
weakens the N--N bond, and helps ammonia leave.

\textbf{11.} The \textit{ortho} carbon has become $sp^3$, with a hydrogen: loss of that proton
restores the benzene ring, a strong driving force.

\textbf{12.} The aniline $\ce{NH2}$ formed attacks the imine carbon, closing the
five-membered ring (a 2-aminoindoline).

\textbf{13.} Ammonia; the formation of the aromatic pyrrole ring of indole drives it.

\textbf{14.} $\ce{C9H12N2 -> C9H9N + NH3}$.

\textbf{15.} \qty{131.0}{g/mol}.

\textbf{16.} $\qty{0.0926}{mol} \times \qty{131.0}{g/mol} = \qty{12.1}{g}$.

\textbf{17.} At C3, the most nucleophilic position of indole, free here; attack there
keeps the benzene ring intact in the intermediate.

\textbf{18.} $\ce{C6H5NHNH-C(=CH2)-CH2CH3}$ (terminal) and $\ce{C6H5NHNH-C(CH3)=CH-CH3}$ (internal).

\textbf{19.} 2-Ethylindole and 2,3-dimethylindole.

\textbf{20.} The internal one is more substituted; strong acids raise the share of the
terminal one.

\textbf{21.} 2,3-Dimethylindole, through the more stable, more substituted
ene-hydrazine.

\textbf{22.} Acetaldehyde condenses with itself and the ene-hydrazine route gives poor
yields; indole is made instead from the hydrazone of pyruvic acid, giving
indole-2-carboxylic acid, which is decarboxylated.

\textbf{23.} A mixture of regioisomers must be separated, which costs yield; a
symmetrical ketone or a route that fixes the regiochemistry is preferred.

\textbf{24.} 10.0 g of phenylhydrazine can give at most \qty{12.1}{g} of 2-methylindole.
\end{solution}
